\documentclass[11pt, a4paper]{article}

% ─── Core packages ───────────────────────────────────────────────────────────
\usepackage[T1]{fontenc}
\usepackage[utf8]{inputenc}
\usepackage{microtype}
\usepackage{lmodern}

% ─── Page geometry ───────────────────────────────────────────────────────────
\usepackage[
  top=2.5cm, bottom=2.5cm,
  left=3cm,  right=3cm,
  marginparwidth=2cm
]{geometry}

% ─── Mathematics ─────────────────────────────────────────────────────────────
\usepackage{amsmath, amssymb, amsthm}
\usepackage{mathtools}

% ─── Tables ──────────────────────────────────────────────────────────────────
\usepackage{booktabs}
\usepackage{tabularx}
\usepackage{array}

% ─── Lists ───────────────────────────────────────────────────────────────────
\usepackage{enumitem}
\setlist[itemize]{noitemsep, topsep=4pt}
\setlist[enumerate]{noitemsep, topsep=4pt}

% ─── Typography ──────────────────────────────────────────────────────────────
\usepackage{setspace}
\setstretch{1.15}
\usepackage{parskip}
\setlength{\parindent}{0pt}
\setlength{\parskip}{6pt}

% ─── Section styling (minimal) ───────────────────────────────────────────────
\usepackage{titlesec}
\titleformat{\section}{\large\bfseries}{\thesection.}{0.5em}{}[\vspace{-4pt}\rule{\linewidth}{0.4pt}]
\titleformat{\subsection}{\normalsize\bfseries}{\thesubsection}{0.5em}{}
\titleformat{\subsubsection}{\normalsize\itshape}{\thesubsubsection}{0.5em}{}
\titlespacing*{\section}{0pt}{14pt}{6pt}
\titlespacing*{\subsection}{0pt}{10pt}{4pt}

% ─── Header / footer ─────────────────────────────────────────────────────────
\usepackage{fancyhdr}
\pagestyle{fancy}
\fancyhf{}
\renewcommand{\headrulewidth}{0.4pt}
\fancyhead[L]{\small\textit{PAPER TITLE}}   % ← replace
\fancyhead[R]{\small\textit{AUTHOR YEAR}}   % ← replace
\fancyfoot[C]{\small\thepage}

% ─── Hyperlinks ──────────────────────────────────────────────────────────────
\usepackage[
  colorlinks=true,
  linkcolor=black,
  citecolor=black,
  urlcolor=black
]{hyperref}

% ─── CJK (uncomment if Chinese content needed) ───────────────────────────────
% \usepackage{xeCJK}

% ─── Custom environments ─────────────────────────────────────────────────────
\newtheorem{finding}{Finding}
\newtheorem{limitation}{Limitation}

% ═════════════════════════════════════════════════════════════════════════════
\begin{document}

% ─── Title block ─────────────────────────────────────────────────────────────
\begin{center}
  {\Large\bfseries PAPER TITLE}\\[6pt]    % ← replace
  {\normalsize AUTHORS}\\[3pt]            % ← replace
  {\small JOURNAL, YEAR}\\[3pt]           % ← replace
  \rule{0.5\linewidth}{0.4pt}\\[3pt]
  {\small\textit{Reading note by} YOUR NAME \quad \today}  % ← replace
\end{center}

\vspace{8pt}

% ─── 1. Core Argument ────────────────────────────────────────────────────────
\section{核心论证}

\textbf{研究问题}

\textbf{论证链}

\textbf{核心贡献}

% ─── 2. Research Design ──────────────────────────────────────────────────────
\section{研究设计}

\subsection{方法概述}

\subsection{识别策略}

\begin{table}[h]
\centering
\small
\begin{tabularx}{\linewidth}{lXl}
\toprule
威胁 & 作者应对 & 充分？ \\
\midrule
 &  &  \\
\bottomrule
\end{tabularx}
\caption{识别策略潜在威胁}
\end{table}

% ─── 3. Data ─────────────────────────────────────────────────────────────────
\section{数据}

\begin{table}[h]
\centering
\small
\begin{tabularx}{\linewidth}{llXl}
\toprule
变量 & 类型 & 构建方式 & 潜在问题 \\
\midrule
 &  &  &  \\
\bottomrule
\end{tabularx}
\caption{关键变量}
\end{table}

% ─── 4. Results ──────────────────────────────────────────────────────────────
\section{核心结果}

\subsection{主要发现}

\begin{table}[h]
\centering
\small
\begin{tabularx}{\linewidth}{lXll}
\toprule
表/图 & 问题 & 核心结果 & 解读 \\
\midrule
 &  &  &  \\
\bottomrule
\end{tabularx}
\caption{主要结果}
\end{table}

\subsection{稳健性检验}

\subsection{机制检验}

\subsection{异质性分析}

% ─── 5. Critical Evaluation ──────────────────────────────────────────────────
\section{批判性评估}

\subsection{方法论问题}

\subsection{外部效度}

\subsection{与文献的张力}

% ─── 6. Idea Seeds ───────────────────────────────────────────────────────────
\section{延伸方向}

\begin{enumerate}
  \item \textbf{Seed 1: TITLE}\\
    \textit{逻辑}: \\
    \textit{RQ}: \\
    \textit{Y / X / 识别}: \\
    \textit{所需数据}: \\
    \textit{难点}:

  \item \textbf{Seed 2: TITLE}\\
    同上结构

  \item \textbf{Seed 3: TITLE}\\
    同上结构
\end{enumerate}

% ─── 7. Replication Notes ────────────────────────────────────────────────────
\section{复现要点}

\end{document}
